% !TeX root = ../papers/nc-p3.tex
\section{核像が曲線となる整八次式の線織正規化と三重母線}
\label{sec:kernel-curve-scroll}
\begin{theorem}\label{thm:kernel-curve-scroll}
四辺有限の標準有限Frobenius窓の整な反対八次曲面 $D$ で、一般核の像 $Z$ が曲線であるとする。
$B$ を $Z$ の正規化、$g$ をその種数、$\delta=\deg Z$ とする。
$D$ の正規化は $B$ 上の射影直線束であり、一般核点に対する主成分の繊維は一本の入力直線である。
母線のPlücker次数は8で、$a=\delta-g-3\ge0$ となる。
余核層は核座標の線束の正規化からの押出しに入り、その商のHilbert多項式は $a(t+1)$ である。
$a=0$ なら両層は等しい。これらに点窓の被約性を要しない。
\end{theorem}
\begin{proof}[余核の双対性を一般繊維で使う]
入力の射影空間で $0\to\mathcal O(-2)^4\xrightarrow{N}\mathcal O^4\to M\to0$ とする。
$M$ は整超曲面 $D$ 上の階数一極大Cohen--Macaulay層で特にtorsion-freeであり、
表示の対称性と双対化から $\mathcal Hom_D(M,\omega_D)=M(-2)$ となる。
核グラフを解消し、その滑らかな曲面を $S$、入力射を $f$、超平面を $H$、核直線の双対を $L$ とする。
四つの核座標から一般点で同型な $M\hookrightarrow f_*L$ を得る。
正規化の導手をあらかじめ特定せず、対称性から自然な包含
\[
      f_*(\omega_S L^{-1})\ \longrightarrow\
      \mathcal Hom_D(M,\omega_D)=M(-2)\ \longrightarrow\ f_*L(-2H)
\]
を使う。第一射は $M$ の評価との積を固有双有理射の標準跡 $f_*\omega_S\to\omega_D$ に送る。
一般点で同型、左端はtorsion-freeなので単射である。
曲線への核射のStein分解を取る。一般連結繊維の近くで $f$ は有限になるよう解消を取れる。
実際入力の固定一点 $p$ に潰れる曲線が核曲線を支配すれば $N(p)$ は一般核全てを消す。
核曲線は非退化（八次式上で消える六次の余因子列は恒等零となる）であり、
$N(p)=0$ は二次式基底に反する。正規化は有限で、その解消の例外曲線も有限本である。
有限な部分で上の合成は分数体での乗法なので正規化の構造層についても線型である。
従って一般繊維 $J$ 上で $2L-K_S-2H$ に非零切断がある。
$LJ=0$ より $2g(J)-2+2HJ\le0$。$HJ$ は正整数なので $g(J)=0,HJ=1$、
各連結繊維は入力直線へ同型に写る。また上の合成は一般繊維上のどの点でも同型である。
従って一般の母線の各点ではこの合成に消滅因子はない。
有限な射の一つの枝で合成が同型なら、その枝の剰余値は左端から独立に指定できる。
複数枝のある点では、全枝で同型なら中間層も押出し全体に一致する。
特に後で使う、一般点で全ての枝が一般母線上にある場合に適用できる。

同じ一般核 $v$ を持つ直線は高々二本である。四つの独立な二次式 $N(x)v$ は
整な八次式の因子となる共通平面や二次曲面を持たず、定理\ref{thm:opposite-column-divisor}の
完全交叉リンケージの次数上界はそのまま適用できる。二本が交わるなら、交点は動く
（固定なら上と同じ理由で $N(p)=0$）。その一般点では二本の異なる正規化枝が同じ核座標を持つ。
交点が他の特別な母線上にもあるとしても、この二枝で先の合成が同型なので、$M$ の評価で
この二つの値を独立に指定できる。実際有限押出しの左端は異なる二つの剰余体に全射し、
合成は各枝で単元倍である。しかし $M$ は四座標で生成され、その二点の核座標は同一、これは矛盾する。
交点の軌跡の一般点で解消が有限でない例外もない。正規化曲面の解消が潰す点は有限である。

二直線がねじれの位置なら、それらを含む二次式は、相補な二つの二変数群の四つの交差積の包で尽きる。
従って $x\mapsto[N(x)v]$ の像はSegre二次曲面に入る。一方一般 $v$ では $y^tv$ は
有限点窓の全台で単元なので、同じ写像は窓上で反対座標 $y$ を与える。
その二次評価の単射性に反する。これでStein分解の底は $Z$ の正規化 $B$ に一致する。

一般の母線を記録する $B\to\operatorname{Gr}(2,4)$ は全点に延長する。
普遍直線族の引戻し $S_0$ は滑らかな射影直線束で、$D$ への双有理射は有限である。
正次元繊維があれば全母線が一点を通るためである。よって $S_0$ が正規化である。
$H$ はその相対超平面で $H^2=8$、すなわちPlücker次数8である。
$L$ は $B$ の次数 $\delta$ の線束の引戻しで、両写像の入射等式は全点へ延びる。
相対超平面の対称冪と曲線のRiemann--Rochより押出しのHilbert多項式は
$(t+1)(4t+\delta+1-g)$、$M$ のものは表示から $4(t+1)^2$。
その差による結論を得る。
\end{proof}

\begin{corollary}\label{cor:kernel-curve-triple-ruling-genus}
上の $a=0$ とし、入力窓は被約、$N$ の階数1軌跡は有限とする。
このとき $g=0$, $\delta=3$ であり、三本の母線の交点を長さ三の繊維として数えるスキームは
被約な二十点である。特に九余因子の核像が滑らかな曲線で滑らかな二次曲面上にあるなら、
この有限階数一・被約窓の場合、その二次数は順序を除き $(1,b)$ に限られる。
\end{corollary}
\begin{proof}
$M=f_*L$ の任意の入力閉点での生成数は、有限射の繊維長に等しく、
$4-\operatorname{rank}N(p)$ である。$N(p)$ は零行列にならないので最大長は3で、
三に達する点は有限個であり、指定二十点は全てその中に入る。
$B$ 上で母線の二平面の商を $F$ とすると、$F$ は次数8の二束、入力の四切断から生成される。
$[p]\in\mathbb P^3$ 上の母線の繊維は対応する切断の零因子である。
$B^{(3)}$ 上の階数六の束 $F^{[3]}$ は、普遍長さ三因子上の $F$ を押し出して取る。
$B^{(3)}\times\mathbb P^3$ で $F^{[3]}\otimes\mathcal O(1)$ の普遍切断の零点は有限である。
どの長さ三繊維もそれに含まれる長さ三因子をただ一つ持つからである。従って長さは
$\int_{B^{(3)}}c_3(F^{[3]})=20-4g$ となる。

計算を明記する。順序付き三乗から次数6で引き戻す。$h_i$ を各点類、$D_i=\sum_{j<i}\Delta_{ij}$ とすれば、
普遍因子を順に加える濾過の三商は $F_i(-D_i)$ である。
従って $d=\deg F$ とし $\prod_i (1+(dh_i-2D_i)t+(D_i^2-dh_iD_i)t^2)$ の三次係数を積分し6で割る。
対角の法線は曲線の接束なので、重複した対角制限は $(2-2g)$ 倍の点類を与え、
結果は $(d-4)((d-2)(d-3)-6g)/6$ である。
被約な二十点が異なる零点を与えるので $g=0$、等長性から全て単純である。
滑らかな二次曲面上の滑らかな $(a_1,b_1)$ 型では $\delta=a_1+b_1$, $g=(a_1-1)(b_1-1)$。
$\delta-g-3=-(a_1-2)(b_1-2)\ge0$ であり、両方が2以上なら片方が2で等号となる。
このとき $g\ge1$ が既に排除した場合である。
\end{proof}
\path{computations/verify_kernel_curve_scroll.py} は順序付き三乗での全対角積を計算し、
任意の $d,g$ の三重点公式を検算する。
